% Gertsenshtein Conversion Formula: Derivation, Conventions, and Literature Comparison
% Source: docs/gertsenshtein_formula.md (migrated to LaTeX)
% Uses macros from preamble.tex

\section{Gertsenshtein Conversion Formula}
\label{sec:gertsenshtein-formula}

\subsection{Summary}

The graviton-photon conversion probability in a uniform magnetic field $\Bzero$ over propagation distance $D$ is:
\begin{equation}
  \Pconv(\text{graviton} \to \text{photon}) = \sin^2\!\left(\frac{\kappa \Bzero D}{2}\right)
  \label{eq:gertsenshtein-P}
\end{equation}
where $\kappa^2 = 16\pi G$ (standard gravitational coupling). Equivalently:
\begin{equation}
  \Pconv = \sin^2\!\left(\sqrt{4\pi G}\, \Bzero\, D\right)
  \label{eq:gertsenshtein-P-alt}
\end{equation}

This is derived from the Lagrangian $\lag = \frac{1}{\kappa^2}R - \frac{1}{4}F^2$ via the TIDAL symbolic pipeline and confirmed numerically (RMS $< 0.015$ across a 40-point $\Bzero$ sweep at $N=1024$, corrected for graviton effective mass). It is independently confirmed by Dandoy, Lella et al.~\cite{dandoy2024constraining}.

\textbf{Palessandro \& Rothman~(2023) quote $P = \sin^2(\sqrt{G}\,\Bzero\,D)$, which is incorrect by a factor of $\sqrt{4\pi} = 2\sqrt{\pi} \approx 3.54$ in the argument.} Domcke \& Garcia-Cely~\cite{domcke2023magnetic} independently criticise P\&R's linearisation procedure.


\subsection{Derivation from the TIDAL Equations}

\subsubsection{Starting Point}

The Lagrangian:
\begin{equation}
  \lag = \frac{1}{\kappa^2}\, R - \frac{1}{4}\, F_{\mu\nu}\, F^{\mu\nu}
  \label{eq:lagrangian}
\end{equation}
with $\kappa^2 = 16\pi G$, background metric $\etamunu = \text{diag}(-1,+1,+1,+1)$, and uniform background magnetic field $\Bzero$ along $x$ generated by $\bar{A}_y = -\Bzero z$.

After:
\begin{itemize}
  \item Perturbation expansion: $g = \eta + \varepsilon h$, $A = \bar{A} + \varepsilon a$
  \item TT gauge on $h$, Lorenz gauge on $a$
  \item Plane-wave reduction ($\ddx = \ddy = 0$, propagation along $z$)
  \item Constraint elimination (traceless $h_4 = -h_7$; transverse $h_6 = h_8 = 0$)
\end{itemize}
the Gertsenshtein channel ($h_7 \leftrightarrow a_2$) satisfies:
\begin{align}
  \ddt^2 h_7 &= \ddz^2 h_7 - \Bzero^2 \kappa^2\, h_7 - \Bzero \kappa^2\, \ddz a_2 \label{eq:eom-h7} \\
  \ddt^2 a_2 &= \ddz^2 a_2 + \Bzero\, \ddz h_7 \label{eq:eom-a2}
\end{align}

\subsubsection{Why the Coupling Is Asymmetric}

The $h_7$ equation has coupling coefficient $\Bzero\kappa^2$ while the $a_2$ equation has $\Bzero$. This arises from the kinetic normalisation step in the pipeline:
\begin{enumerate}
  \item \textbf{VarD} (EulerLagrange.wl) computes $\delta\lag/\delta h = 0$, giving: $(1/\kappa^2)\,\partial^2 h + \text{coupling} = 0$
  \item \textbf{ExtractLHSCoefficient} (ExportJSON.wl) identifies $\text{lhsCoeff} = 1/\kappa^2$
  \item \textbf{LHS normalisation} (ExportJSON.wl): $\text{rhs} = -\text{rhs}/\text{lhsCoeff}$ divides by $1/\kappa^2$
\end{enumerate}
This multiplies the coupling by $\kappa^2$: the raw coupling $\Bzero$ becomes $\Bzero\kappa^2$ in the $h$ equation. The $a$ equation has $\text{lhsCoeff} = 1$ (Maxwell kinetic term), so its coupling $\Bzero$ is unchanged.

\subsubsection{Eigenmode Analysis}

For a plane-wave mode $e^{i(kz - \omega t)}$, substitute $\ddz \to ik$ into \cref{eq:eom-h7,eq:eom-a2}:
\begin{align}
  \omega^2 H &= (k^2 + \Bzero^2\kappa^2)\,H + ik\,\Bzero\kappa^2\, A \\
  \omega^2 A &= k^2\, A - ik\,\Bzero\, H
\end{align}
giving the mixing matrix:
\begin{equation}
  M = \begin{pmatrix}
    k^2 + \Bzero^2\kappa^2 & ik\,\Bzero\kappa^2 \\
    -ik\,\Bzero & k^2
  \end{pmatrix}
  \label{eq:mixing-matrix}
\end{equation}

\textbf{Eigenvalues:} $\det(M) = k^4$, $\text{Tr}(M) = 2k^2 + \Bzero^2\kappa^2$, giving:
\begin{equation}
  \lambda_\pm = \frac{2k^2 + \Bzero^2\kappa^2 \pm \Bzero\kappa\sqrt{4k^2 + \Bzero^2\kappa^2}}{2}
\end{equation}

\textbf{For $k \gg \Bzero\kappa$} (massless dispersion limit):
\begin{equation}
  \omega_\pm \approx k \pm \frac{\Bzero\kappa}{2}
\end{equation}

\textbf{Beat frequency:} $\Delta\omega = \omega_+ - \omega_- = \Bzero\kappa$. The conversion probability is therefore:
\begin{equation}
  \boxed{\Pconv(t) = \sin^2\!\left(\frac{\Delta\omega\, t}{2}\right) = \sin^2\!\left(\frac{\kappa\,\Bzero\, t}{2}\right)}
  \label{eq:P-final}
\end{equation}

Note: the beat frequency $\Bzero\kappa$ is the geometric mean of the two asymmetric couplings: $\sqrt{\Bzero\kappa^2 \cdot \Bzero} = \Bzero\kappa$. This is a general result for coupled oscillators with asymmetric coupling strengths.


\subsection{Convention Comparison Across the Literature}
\label{sec:conventions}

\begin{table}[htbp]
\centering
\caption{Conversion formula by convention. All three rows are algebraically equivalent except P\&R, which is incorrect by $\sqrt{4\pi}$.}
\label{tab:conventions}
\begin{tabular}{@{}lllll@{}}
\toprule
Source & Metric pert.\ & Gravity kinetic & Coupling & Vacuum $P$ \\
\midrule
\textbf{TIDAL} & $g = \eta + h$ & $(1/\kappa^2)(\partial h)^2$ & $\kappa\Bzero/2$ & $\sin^2(\kappa\Bzero D/2)$ \\
\textbf{Dandoy et al.}~\cite{dandoy2024constraining} & $g = \eta + \tfrac{2}{M_P}h$ & $\tfrac{1}{2}(\partial h_c)^2$ & $\Bzero/(\sqrt{2}\,M_P)$ & $\sin^2(\sqrt{4\pi G}\,\Bzero D)$ \\
\textbf{P\&R (2023)} & $g = \eta + h$ & $(1/16\pi G)(\partial h)^2$ & $\sqrt{G}\,\Bzero$ & $\sin^2(\sqrt{G}\,\Bzero D)$ \\
\bottomrule
\end{tabular}
\end{table}

\subsubsection{Algebraic Equivalence: TIDAL = Dandoy/Lella}

TIDAL uses $\kappa^2 = 16\pi G$:
\begin{equation}
  \frac{\kappa}{2} = \frac{\sqrt{16\pi G}}{2} = \sqrt{4\pi G}
\end{equation}
Dandoy/Lella use the reduced Planck mass $M_P = 1/\sqrt{8\pi G}$:
\begin{equation}
  \frac{\Bzero}{\sqrt{2}\,M_P} = \Bzero\,\frac{\sqrt{8\pi G}}{\sqrt{2}} = \Bzero\,\sqrt{4\pi G}
\end{equation}
These are identical: $\kappa/2 = \sqrt{4\pi G} = \Bzero/(\sqrt{2}\,M_P)$~\checkmark

\subsubsection{P\&R's Error}

P\&R use $g = \eta + h$ with $\lag = -(1/16\pi G)\,(\partial h)^2$. Their graviton kinetic term is non-canonical. The correct mixing element is $\sqrt{4\pi G}\,\Bzero$; they obtain $\sqrt{G}\,\Bzero$, missing a factor of $\sqrt{4\pi} \approx 3.54$.

The error ratio in the weak-field limit:
\begin{equation}
  \frac{P_\text{TIDAL}}{P_\text{PR}} = \frac{\sin^2(\sqrt{4\pi G}\,\Bzero D)}{\sin^2(\sqrt{G}\,\Bzero D)} \approx \frac{4\pi G\,\Bzero^2 D^2}{G\,\Bzero^2 D^2} = 4\pi
\end{equation}


\subsection{Graviton Effective Mass and $P_\text{max}$ Correction}
\label{sec:mass-correction}

\subsubsection{Physical Origin}

The $h_7$ equation contains a mass term $-\kappa^2 \Bzero^2\, h_7$, acting as an effective mass squared $m_\text{eff}^2 = \kappa^2 \Bzero^2$ for the graviton, arising from the background EM stress-energy. The photon $a_2$ is massless.

\subsubsection{Effect on Energy-Based Conversion Measurement}

TIDAL measures conversion as an energy ratio $P(t) = E_\text{target}(t) / E_\text{source}(0)$. For initial condition $h_7 = A\cos(kz)$:
\begin{align}
  E_\text{source}(0) &\propto \frac{k^2 + \kappa^2\Bzero^2}{2\kappa^2}\, A^2 L \quad\text{(graviton with mass)} \\
  E_\text{target}(t_\text{peak}) &\propto \frac{k^2}{2\kappa^2}\, A^2 L \quad\text{(massless photon)}
\end{align}
Therefore:
\begin{equation}
  P_\text{max} = \frac{k^2}{k^2 + \kappa^2 \Bzero^2}
  \label{eq:Pmax}
\end{equation}
The full corrected formula:
\begin{equation}
  \boxed{P_\text{corrected}(\Bzero, t) = \sin^2\!\left(\frac{\kappa\,\Bzero\, t}{2}\right) \times \frac{k^2}{k^2 + \kappa^2\,\Bzero^2}}
  \label{eq:P-corrected}
\end{equation}

\subsubsection{Numerical Verification}

Targeted simulations ($k = 2.0106$, $\kappa = 1$, $t_\text{end} = 50$, domain $[0,100]$):

\begin{table}[htbp]
\centering
\caption{$P_\text{max}$ comparison: TIDAL simulation vs.\ corrected formula \cref{eq:Pmax}.}
\label{tab:Pmax-verification}
\begin{tabular}{@{}cccc@{}}
\toprule
$\Bzero$ & $P_\text{max}$ (TIDAL) & $k^2/(k^2+\kappa^2\Bzero^2)$ & Difference \\
\midrule
0.0628 & 0.9984 & 0.9990 & $-0.0006$ \\
0.1000 & 0.9982 & 0.9975 & $+0.0007$ \\
0.1885 & 0.9948 & 0.9913 & $+0.0035$ \\
0.2000 & 0.9941 & 0.9902 & $+0.0039$ \\
\bottomrule
\end{tabular}
\end{table}

All measurements agree with the corrected formula to $< 0.4\%$.


\subsection{Numerical Dispersion}
\label{sec:numerical-dispersion}

The Rabi period measured from TIDAL is systematically longer than $\pi/(\kappa\Bzero)$, with the discrepancy scaling as $O(k\Delta x)^2$:

\begin{table}[htbp]
\centering
\caption{Numerical dispersion: effective Rabi frequency vs.\ grid resolution.}
\label{tab:dispersion}
\begin{tabular}{@{}cccc@{}}
\toprule
$N$ (grid points) & $k\Delta x$ & $\Omega_\text{eff}/\Omega_\text{theory}$ & Error \\
\midrule
512  & 0.393 & 0.9802 & 1.98\% \\
1024 & 0.196 & 0.9943 & 0.57\% \\
2048 & 0.098 & $\approx 0.999$ & $\sim 0.15\%$ \\
\bottomrule
\end{tabular}
\end{table}

The $4\times$ error reduction from doubling $N$ confirms the $O(\Delta x^2)$ scaling of the central-difference spatial operators. \textbf{Practical implication:} use $N \geq 1024$ grid points for quantitative $\Bzero$ sweep comparisons.


\subsection{Detuned Generalisation}

The formula generalises to the detuned case (plasma effects, massive graviton):
\begin{equation}
  P = \left(\frac{\mu}{\omega_m}\right)^2 \sin^2(\omega_m\, D),
  \qquad \mu = \frac{\kappa\Bzero}{2}, \quad
  \omega_m = \sqrt{\Delta^2 + \mu^2}, \quad
  \Delta = \frac{m_\gamma^2 - m_g^2}{2\omega}
  \label{eq:detuned}
\end{equation}

\textbf{Weak-field limit:} $P \approx \kappa^2 \Bzero^2 D^2 / 4 = 4\pi G\, \Bzero^2 D^2$ (not $G\,\Bzero^2 D^2$ as P\&R claim).


\subsection{Literature Notes}
\label{sec:literature-notes}

\begin{description}
  \item[Gertsenshtein (1962)~\cite{gertsenshtein1962wave}] Original prediction. Weak-field limit: $P \approx (\kappa B L/2)^2$, consistent with $\sin^2 \approx x^2$ for small argument.
  \item[Raffelt \& Stodolsky (1988)] Established the $2\times 2$ mixing-matrix formalism by analogy with neutrino oscillations. Introduced MSW resonance for graviton-photon mixing.
  \item[Domcke \& Garcia-Cely (2023)~\cite{domcke2023magnetic}] Explicitly criticises P\&R's treatment of the metric in the field strength tensor.
  \item[Dandoy et al.\ (2024)~\cite{dandoy2024constraining}] Full 4-component mixing (two GW polarisations $\times$ two photon polarisations). \textbf{Confirms TIDAL's formula:} $P = \sin^2(\sqrt{4\pi G}\,\Bzero D)$.
  \item[Domcke et al.\ (2025)~\cite{domcke2025scattering}] GW scattering on localised magnetic fields. Extension target for TIDAL.
\end{description}


\subsection{Implications}

\begin{enumerate}
  \item \textbf{Our formula is correct.} The TIDAL pipeline (xAct $\to$ VarD $\to$ ExportJSON normalisation) correctly derives the coupled equations.
  \item \textbf{P\&R is not a reliable validation target.} The 2023 and 2024 papers both contain the same normalisation error. Dandoy et al.~\cite{dandoy2024constraining} should be used as the primary cross-check.
  \item \textbf{The formula generalises cleanly} to the detuned case via \cref{eq:detuned}.
\end{enumerate}
